\clearpage
\section*{September 24, 2026: Same-owner companion filings join one parent}
\textit{Agent-drafted research entry.}

The production linker joined filings only through a shared lot, lot history,
an explicit project reference or exact lot adjacency. A developer who files
two buildings on separate lots across the street therefore appeared as two
parents. Jacob asked for one symmetric rule instead of case-by-case research,
tested it in \path{audit_companion_rules}, and asked that it enter parent
construction for both periods.

\paragraph{The rule.} Two filings are companions when the owner business name
or the owner person named on their own applications matches, they are within
150 metres or on the same tax block within 200 metres, and they were filed
within 365 days. DOB NOW jobs take the owner from DOB NOW; older BIS jobs
(historical filings through 2020) from a new dated capture of BIS New Building
applications, \path{fetch_dob_bis_job_filings}. Placeholders such as
\texttt{PR} or \texttt{NOT APPLICABLE}, public agencies, and people who sign as
owner for an agency do not identify a developer. Coordinates are Housing
Database coordinates with DOB fallback. The links are applied after the
existing ones, under the same 365-day limit from a parent's first to last
filing. Manually rejected pairs cannot be joined, directly or through a third
filing.

The 200-metre limit on same-block links came from one case. Beitel's 261--315
Grand Concourse and 120 East 144th Street share a very large block and an
owner, but public records describe them as separate projects.

\paragraph{Validation.} The September 23 hand review covers 111 links. Of the
40 random historical links under the rule, 34 are clear companions, five
likely companions and one separate. Of the 31 recent links, 18 are clear, ten
likely and three uncertain. Links at 150--250 metres are weaker (9 clear, 4
likely, 6 uncertain, 1 separate), which is why the distance stops at 150.
Every reviewed link within 60 metres is a companion. Links 60--150 metres apart
filed months apart are the mixed group.

\paragraph{What changed.} The rule adds 1,058 historical and 106 recent links
that no existing rule found, and parent membership falls from 6,828 to 6,440
historical and from 1,943 to 1,873 recent parents. Filings and units are
unchanged. Among A/B rental parents with at least 50 units:

\begin{center}
\small
\begin{tabular}{lrrrr}
\hline
& \multicolumn{2}{c}{Historical} & \multicolumn{2}{c}{Post-policy} \\
& Before & After & Before & After \\
\hline
Parents & 606 & 599 & 327 & 319 \\
Multi-building share & 8.4\% & 13.9\% & 19.3\% & 23.5\% \\
Three or more buildings & 2.5\% & 3.5\% & 4.9\% & 6.3\% \\
Exact-99 share & 0.5\% & 0.5\% & 11.6\% & 11.6\% \\
Two 99-unit buildings (99+99) & 0 & 0 & 9 & 9 \\
\hline
\end{tabular}
\end{center}

Splitting is more common in both periods than the old links showed, by about
five points historically and four after the policy. The reweighted
comparison barely moves. The weighting sample is 588 historical and 304 recent
parents (595 and 310 before), the effective sample size 519, the excess at 99
11.3 points (11.4 before; bootstrap interval 7.6--15.2), the excess at 198 3.0
points, and the 100--149 deficit 7.6 points (7.3 before; 3.1--11.6).

Joining filings changes which filing anchors a parent. Five historical and
four recent constituents joined an earlier filing's parent anchored before the
comparison window (in 2018 and 2024) and so left the panels; one is 545
Vanderbilt Avenue, 264 units. Two reviewed land decisions in
\path{site_lot_decisions.csv} were keyed to parents that now include another
filing, and were moved to the new parent with that filing's lot added: CAMBA's
filings 210179965 and 210179974, and Grun's 1973 Crotona Avenue and 623 East
178th Street.

The remaining 13 historical and 2 recent rule matches are not applied because
the joined parent would span more than 365 days. Re-forming the production
parents under the rule in the audit changes none of them.

\paragraph{Open question.} The rule cannot distinguish a developer's companion
buildings from its separate nearby projects. The Grun link matters for the
notch: 623 East 178th Street was a 99-unit parent and is now joined with 1973
Crotona Avenue, 44 units and 126 metres away, filed 75 days earlier, so the
recent 99-unit weighting-sample parents fall from 37 to 36. Public reporting
describes 623 East 178th Street as a separate Grun project; the review now
marks the link uncertain. Six historical links in the weighting sample rest on
the owner person alone, with different LLCs. Two look like companions (1731
and 1745 West Farms Road; 4829 and 4831 White Plains Road), one is uncertain
(280 Meeker and 680 Lorimer Street), one appears separate (389 Weirfield and
1504 Jefferson Street), and two were not researched. Whether the rule should
narrow, for example by requiring the same architect for person-only links or a
shorter window beyond 60 metres, is Jacob's decision.

\textit{Reproduction note.} The links are built by
\path{construct_owner_proximity_links.R} in \path{construct_parent_cohorts};
root \texttt{make data}. Root \texttt{make companion-rules} scores the
alternative rules and simulates their parents against the production panel.

\subsection*{September 24 follow-up: the rule stays, with a stricter check}

Jacob kept the rule as adopted and asked for a stricter version as a
sensitivity check. The strict version drops the links supported only by a
shared owner unless the filings are within 60 metres or filed within 30 days,
where every reviewed link was a companion. It removes 285 historical and 24
recent links and splits 9 historical and 6 recent weighting-sample parents,
among them Grun's 1973 Crotona Avenue and 623 East 178th Street and the 280
Meeker Avenue and 680 Lorimer Street pair. Each piece keeps its parent's site
traits and classification, pieces below 50 units leave the sample, and the
historical weights are recalibrated.

\begin{center}
\small
\begin{tabular}{lrr}
\hline
& Adopted rule & Strict owner links \\
\hline
Historical / post parents & 588 / 304 & 593 / 303 \\
Excess at 99 & 11.3 & 11.7 \\
Excess at 198 & 3.0 & 3.0 \\
Deficit 100--149 & 7.6 & 7.8 \\
Post exact-99 parents & 36 & 37 \\
Post 99+99 parents & 9 & 9 \\
Historical multi-building share (weighted) & 10.7\% & 9.6\% \\
Post multi-building share & 20.7\% & 18.8\% \\
\hline
\end{tabular}
\end{center}

The bunching comparison does not depend on the links in dispute: the excess at
99 moves by 0.4 points, well inside its bootstrap interval, and the 99+99 count
is unchanged. The multi-building shares fall by one to two points in both
periods. The check does not recover the nine filings that production joined
to parents anchored before the comparison window.

\textit{Reproduction note.} \path{strict_rule_sensitivity.R} in
\path{audit_companion_rules}; root \texttt{make companion-rules}. The adopted
row reproduces the benchmark weights' effective sample size and moments.

\subsection*{September 24: community district in the panel}

The parent panel now carries \texttt{community\_district}, the MapPLUTO
district of the parent's lots in the same release that supplies its land and
zoning, so both periods are measured the same way. No weighting-sample parent
lacks one; three in each period span two districts (\texttt{Mixed}). The
sample covers 55 historical and 54 recent districts. Five districts have
parents in only one period, the largest being Brooklyn 7 (Sunset Park), with
seven historical parents and none after the policy, so a calibration on
districts would need to pool small ones. Nothing else in the panels changed.

\subsection*{September 24: weights match zoning and borough}

Jacob decided that the historical weights should not match on lot area. A
parent's land is part of the developer's choice, through assembly, so matching
on it could remove part of the policy response. The benchmark now matches
residential FAR and borough. Built FAR is also left out: it is existing floor
divided by the parent's own lot area. Lot area returns as a robustness check,
alone and in the earlier four-trait match. District was not added, because
the post-policy mix of districts may itself respond to the policy.

\begin{center}
\small
\begin{tabular}{lrrrrr}
\hline
Match & ESS & Excess at 99 & Deficit 100--149 & Hist.\ mean & Hist.\ multi \\
\hline
None & 588 & 11.3 & 8.1 & 176 & 12.2\% \\
Zoning and borough (adopted) & 561 & 11.3 & 7.6 & 171 & 12.5\% \\
Adding lot area & 520 & 11.3 & 7.6 & 146 & 10.7\% \\
Previous (lot area, built FAR) & 519 & 11.3 & 7.6 & 146 & 10.7\% \\
\hline
\end{tabular}
\end{center}

The post-period mean is 140 units and its multi-building share 20.7 percent.
Under the adopted weights the bootstrap intervals are 7.5--15.2 points for the
excess at 99 and 3.4--11.5 for the 100--149 deficit; the stricter
companion rule gives 11.7 and 7.8.

The notch comparison is the same under every match. Lot area matters for
size. Recent parents sit on smaller lots, so matching on lot area lowers the
historical mean from 171 to 146 units. Without it, historical parents are about
30 units larger than recent ones; with it, about 7. How much of that gap is a
response to the policy (developers assembling less land) rather than a change
in available sites is the question the land choice leaves open. It matters for
any estimate of housing forgone, not for the bunching at 99. Under the adopted
weights the land-measurement scenarios change nothing except through the
parents they drop; they now bear only on the lot-area checks.

\textit{Reproduction note.} The formulas are in
\path{tasks/shared/code/scale_shape_helpers.R}; root \texttt{make reweighting}
for the benchmark and \texttt{make land-sensitivity} for the robustness
matches (column \texttt{match}). The exploratory pure-notch pilot has not been
rerun.
